% B4横・2段組の演習プリント解説テンプレート
% 対応する演習プリント: 繰数N_単元名.tex
\documentclass[paper=b4j,landscape,twocolumn,fleqn]{jlreq}
\special{papersize=\the\paperwidth,\the\paperheight}
\setlength{\columnseprule}{0.4pt}
\pagestyle{empty}

\usepackage[margin=8truemm]{geometry}
\usepackage[dvipdfmx]{graphicx}
\usepackage{amsmath,amssymb}
\usepackage{tikz}
\usepackage{pgfplots}
\pgfplotsset{compat=1.18}
\usetikzlibrary{arrows.meta,calc,angles,quotes}

\ModifyHeading{section}{lines=1,before_space=6pt,after_space=0pt}
\setlength{\parskip}{0pt}
\setlength{\jot}{3pt}

%==========================================================
% 解説ブロック用コマンド
%   \kangaekata{考え方の説明}
%   \tochushiki{途中式の align* 環境の中身}
%==========================================================
\newcommand{\kangaekata}[1]{\vspace{2pt}\noindent\textbf{考え方}\quad #1\par\vspace{2pt}}
\newcommand{\tochushiki}[1]{%
  \noindent\textbf{途中式（代表例）}
  \begin{align*}
    #1
  \end{align*}
}

\begin{document}
\normalsize

%==========================================================
\section{大問1のタイトル \quad 解答}
%==========================================================
\begin{align*}
&1.\  \text{答え}\\[0.45em]
&2.\  \text{答え}\\[0.45em]
&3.\  \text{答え}
\end{align*}

\kangaekata{この問題を解くための着目点・方針を書く。}

\tochushiki{
  &1.\ \text{式変形の手順}\\
  &\quad = \text{途中} = \text{答え}
}

% 図説が必要な場合は以下のブロックを使う
%\noindent\textbf{図説}
%\begin{center}
%\begin{tikzpicture}[scale=0.9]
%  % tikz コードをここに記述
%\end{tikzpicture}
%\end{center}

\newpage

%==========================================================
\section{大問2のタイトル \quad 解答}
%==========================================================
\begin{align*}
&1.\  \text{答え}\\[0.45em]
&2.\  \text{答え}\\[0.45em]
&3.\  \text{答え}
\end{align*}

\kangaekata{この問題を解くための着目点・方針を書く。}

\tochushiki{
  &1.\ \text{式変形の手順}\\
  &\quad = \text{途中} = \text{答え}
}

%\noindent\textbf{図説}
%\begin{center}
%\begin{tikzpicture}[scale=0.9]
%  % tikz コードをここに記述
%\end{tikzpicture}
%\end{center}

\end{document}
